\section*{v1.4.0 Derivation and Appendix Integration Core}
\addcontentsline{toc}{section}{v1.4.0 Derivation and Appendix Integration Core}

\paragraph{Normalization note.} In Version~\docversion, this block is read as a \emph{mixed historical derivation layer}: its closure-ladder and micro-to-macro passages belong genealogically to $\mathfrak G$ and $\mathfrak P$, while its probe-reduction and appendix-coupling material remains operational support for the normalized architecture. It is therefore preserved for proof-memory and appendix linkage, but not as a second normative origin.


The specific task of v1.4.0 is to turn the earlier consolidation and
lemma-development layers into a more tightly chained proof surface. This means
that the argument is no longer presented only as a list of strategic claims.
Instead, the intermediate dependencies are written so that later appendix items
can be cited as reusable proof modules. The rule of this stage remains simple:
\emph{no shortening, no creative overwrite, only denser local chaining and
clearer main-text/appendix interlock.}

\subsection*{1. Closure ladder as a reusable derivation chain}
The closure architecture is now fixed as a ladder rather than a single jump.
The intended reading is
\[
\begin{aligned}
&\text{local triolic unit} \Longrightarrow \text{orthogonal stabilization}\\
&\Longrightarrow \text{compensated identification} \Longrightarrow \text{closure class}\\
&\Longrightarrow \text{effective projection}.
\end{aligned}
\]
The advantage of this ladder is methodological. It prevents the later sections
from collapsing several distinct proof moves into one rhetorical transition.
Each arrow can now be isolated, tested, and, where needed, exported into the
appendix as a technical lemma or shortcut.

\begin{lemma}[Orthogonal stabilization lemma]
\label{lem:orthogonal-stabilization-lemma}
Whenever one member of the triolic organization remains orthogonal to the other
 two, the local structural body is prevented from degenerating into a freely
rotating two-sector ambiguity. The third direction therefore acts as a
stabilizing anchor for all later compensator and closure statements.
\end{lemma}

\begin{proof}
Without the orthogonal anchor, the first two directions would admit a family of
relabelings that preserve local appearance while destroying global structural
identity. The orthogonal third component blocks this ambiguity by fixing a
non-coincident direction that all admissible couplings must respect.
\end{proof}

\begin{lemma}[Compensated identification lemma]
\label{lem:compensated-identification-lemma}
If the same structural body is observed through shifted or dual frames, then a
compensating map is forced in order to preserve common identity across those
frames. In the present framework this map is read as the Heisenberg
compensator.
\end{lemma}

\begin{proof}
Frame-shifted representatives are not independent ontological copies. They are
multiple presentations of one common body. A balancing operation is therefore
necessary to identify them as one object rather than a divergent family. That
operation is exactly what HC contributes at the structural level.
\end{proof}

\begin{lemma}[Closure-class lemma]
\label{lem:closure-class-lemma}
Once orthogonal stabilization and compensated identification are both active,
the admissible states assemble into a closure class: transition is allowed,
but leakage into an incompatible external regime is not.
\end{lemma}

\begin{proof}
The orthogonal anchor limits admissible couplings, while HC neutralizes those
frame-induced mismatches that would otherwise split the common body. The
surviving states therefore form a self-consistent class under the allowed
moves of the theory.
\end{proof}

\subsection*{2. Micro-to-macro chain tightened}
The earlier versions already stated that the macroscopic projection must be read
from the microscopic structural unit through HC and closure. In v1.4.0,
this chain is stated more sharply:
\[
\begin{aligned}
&\text{micro structural unit} \Longrightarrow \text{local compensated coupling}\\
&\Longrightarrow \text{sector stability} \Longrightarrow \text{macroscopic effective law}.
\end{aligned}
\]
No macroscopic term is treated as primitive if it can be located lower in this
chain. This applies in particular to geometric large-scale terms, matter/gauge
labels, and all observer-frame residues.

\begin{proposition}[Micro-to-macro closure principle]
\label{prop:micro-to-macro}
Every effective macroscopic equation used in the Sphaera-Cascade line must be
read as a closure-compatible projection of the local triolic body rather than
as an independent starting law.
\end{proposition}

\begin{proof}
The ground condition, orthogonality requirement, and HC mechanism already fix
what counts as an admissible local evolution. A macroscopic law that failed to
arise as a projection of that admissible class would describe a sector not
licensed by the common body. Hence effective macroscopic forms are subordinate
projections of the closure class.
\end{proof}

\paragraph{Interpretive consequence.}
This principle strengthens the earlier interpretive discipline. Einstein-like,
Dirac-like, and QFT-like forms remain useful projective readings, but their
status is now tethered to a single chain of legitimacy: they are admissible
only insofar as they compress the closure-compatible microstructure.

\subsection*{3. Probe-reduction and guided search principle}
The probe sector is now equipped with a more explicit reduction principle. Two
probes approaching one another do not explore an arbitrary ambient continuum;
they reduce the candidate space by repeated closure tests.

\begin{lemma}[Probe-reduction lemma]
\label{lem:probe-reduction-lemma}
Let each probe carry its own triolic organization, and let the third group be
orthogonal to both. Then the admissible interaction region is the subset of the
naive search space that survives orthogonality, compensator compatibility, and
closure tests simultaneously.
\end{lemma}

\begin{proof}
Each probe contributes an internal triolic admissibility condition. The shared
orthogonal third group cuts away probe pairings that cannot be embedded in one
common constrained geometry. HC removes residual representational drift, and
closure excludes the remaining leakage channels. What survives is exactly the
reduced search region relevant to structural read-off.
\end{proof}

\paragraph{Tool interlock.}
This is the point where the normalized tools become operational rather than
merely classificatory. Kasisky is attached to frequency-pattern selection.
Hash/rainbow devices compress routing across candidate branches. Markov control
organizes transitions inside the still-admissible region, and Newton
accelerates convergence once a branch has already been selected. The appendix
can therefore store these devices as modular search operators without severing
them from the main proof line.

\subsection*{4. Appendix coupling rules}
The appendix is no longer a passive storage area. In v1.4.0 it is given a
precise function: each appendix block must either (i) justify a compressed step
in the main text, (ii) codify a reusable tool, or (iii) collect a repeated
technical transformation that would otherwise bloat the main line.

The intended appendix map is therefore:
\begin{enumerate}[leftmargin=2em]
  \item \textbf{Symbol and sector tables} collect notation and prevent local
  restatement overhead.
  \item \textbf{Structural tool records} store normalized rechen-tricks in the
  six-field format already fixed earlier.
  \item \textbf{Technical lemmas} receive derivations whose content is
  load-bearing but too local or repetitive for the main line.
  \item \textbf{Projection tables and read-off schemes} record the disciplined
  passage from structural objects to effective sector forms.
\end{enumerate}

\begin{proposition}[Appendix transfer principle]
\label{prop:appendix-transfer}
A step may be moved from the main text to the appendix only if its logical role
is preserved through an explicit citation path. Appendix displacement is
therefore allowed for expansion, but never for concealment.
\end{proposition}

\begin{proof}
The project rule forbids silent loss of proof substance. Therefore, every step
outsourced to the appendix must remain recoverable as a named lemma, tool, or
transformation block. Once that citation path exists, the main text may use the
compressed reference without reducing logical transparency.
\end{proof}

\subsection*{5. v1.4.0 readiness criteria}
This release should be regarded as complete in working form when the following
conditions hold simultaneously:
\begin{enumerate}[leftmargin=2em]
  \item the version stamp and DOI are updated to v1.4.0;
  \item the closure ladder is stated as an explicit reusable chain;
  \item at least one new family of intermediate lemmas is present;
  \item the micro-to-macro principle is written as a formal proposition;
  \item the probe sector is tied directly to the normalized tool apparatus;
  \item the appendix is assigned explicit transfer rules instead of being a
  passive dump area.
\end{enumerate}

These conditions are now satisfied in the present monolith. Relative to v1.3.0,
v1.4.0 therefore marks the point where the document stops being only a
consolidated development line and becomes a more disciplined proof-and-appendix
system.


